\begin{center}\begin{tikzpicture}[font=\small,>=stealth,line join=round]
\draw[gray!40,rounded corners=3pt] (0,0) rectangle (14.2,8.3);\node at (7.1,7.700000000000001) {偶函数的右半边递增：最小值在中间};
\begin{scope}[shift={(4.1,2.0951612903225807)},x=0.5263157894736842cm,y=0.14838709677419354cm]\draw[->,gray] (-5.7,0)--(5.7,0) node[right]{$x$};\draw[->,gray] (0,-3)--(0,28) node[above]{$y$};\draw[blue,thick,solid,domain=-5:5,samples=180] plot(\x,{(\x)^2});\fill[orange!85!black] (-4,16) circle(2.2pt);\fill[orange!85!black] (-2,4) circle(2.2pt);\fill[orange!85!black] (0,0) circle(2.2pt);\fill[orange!85!black] (2,4) circle(2.2pt);\end{scope}\node[anchor=west,align=left,text width=5.0cm] at (8.4,4.0) {用 $f(x)=x^2$ 示范\\[9pt]定义域 $[-5,5]$\\[9pt]$f(-4)=16,\ f(2)=4$\\[9pt]$f(0)=0$\\[9pt]$f(-2)=f(2)$\\[9pt]一般结论：$f(0)$ 最小};
\node[text width=12.6cm,align=center] at (7.1,.65) {最小值所在位置由对称和右半边递增决定；具体最小值不一定是0。};
\end{tikzpicture}\end{center}
